%% kapa: exam=1 round=34 year=2023 period=1 form=A subject=economics subject_ko=경제학원론 seq=39 number=79
%% source: Q-Net 감정평가사 기출문제 (cst00309) artlSeq=5211789 file=raw/1008/5211789/2023년도 제34회 감정평가사 1차시험 1교시.pdf
%% conversion: kordoc 4.24.51 -> tools/gen1.py (auto); review=proofread-sonnet; flags=none
\begin{question}{79}
보호무역을 옹호하는 주장의 근거가 아닌 것은?
\begin{choices}
  \item 자유무역으로 분업이 강력하게 진행되면 국가 안전에 대한 우려가 발생할 수 있다.
  \item 관세를 부과하면 경제적 순손실(deadweight loss)이 발생한다.
  \item 환경오염도피처가 된 거래상대국으로부터 유해한 물질이 자유무역으로 인해 수입될 가능성이 높다.
  \item 정부가 신생 산업을 선진국으로부터 보호해서 육성해야 한다.
  \item 자유무역은 국내 미숙련근로자의 임금에 부정적 영향을 줄 수 있다.
\end{choices}
\end{question}
\answer{2}
